\chapter{Probability and Information for Cryptography}

Cryptography is the art of making the right things unpredictable. A key an attacker cannot guess, a ciphertext that looks like noise, a nonce that never repeats -- every one of these is a statement about probability. If you cannot reason about how likely Eve is to guess a secret, you cannot say a single useful thing about whether a scheme is secure.

This chapter builds the probability toolkit you will lean on for the rest of the book. We start with the bare essentials -- sample spaces, uniform keys, independence -- using tiny concrete numbers you can check by hand. Then we get to the two ideas that do the heaviest lifting in real cryptography: \emph{entropy} (how much unpredictability a secret carries on average) and its paranoid cousin \emph{min-entropy} (how predictable the secret is in the single worst case). We work through the birthday paradox, which is the reason a $256$-bit hash gives you only $128$ bits of collision resistance, and we close with \emph{statistical distance}, the formal way to say ``these two distributions look the same to anyone watching.''

None of this requires measure theory. Everything lives on finite sets, where probability is just counting done carefully.

\chapterline

\section{Why probability is the native language of secrecy}

Picture the standard cast. Alice wants to send Bob a message. Eve sits on the wire and sees everything that crosses it. The central question is: can Eve learn something she should not?

Eve almost never has zero information and almost never has perfect information. She has \emph{partial} information, and partial information is exactly what probability measures. When we say ``a $128$-bit key is secure,'' we are not claiming Eve \emph{cannot} guess it -- she can, with probability $2^{-128}$ per guess. We are saying that probability is so small that even a planet-sized computer running for the age of the universe will not stumble onto it. Security is quantitative, and the quantity is a probability.

\begin{intuitionbox}
A secret is only as good as it is unpredictable. ``Unpredictable'' is not a vibe; it is a number. A coin you have already seen land heads has zero unpredictability. A freshly flipped fair coin has exactly one bit of it. The entire job of a key-generation routine is to manufacture unpredictability and hand it to you in a usable form.
\end{intuitionbox}

\subsection{Sample spaces, events, and uniform distributions}

A finite probability space is a finite set $\Omega$ (the \emph{sample space}) of possible outcomes, together with a function $\Pr$ assigning each outcome $\omega \in \Omega$ a probability $\Pr[\omega] \ge 0$, such that $\sum_{\omega \in \Omega} \Pr[\omega] = 1$. An \emph{event} $A \subseteq \Omega$ is any set of outcomes, and $\Pr[A] = \sum_{\omega \in A} \Pr[\omega]$.

The distribution we use more than any other is the \emph{uniform} one. If $\abs{\Omega} = N$, the uniform distribution puts $\Pr[\omega] = 1/N$ on every outcome. For an event $A$,
\[
  \Pr[A] = \frac{\abs{A}}{\abs{\Omega}} = \frac{\text{favorable outcomes}}{\text{total outcomes}}.
\]
This is the ``count and divide'' rule you already know, and it is the bedrock of key generation: a uniformly random $k$-bit key is a uniform draw from $\Omega = \set{0,1}^k$, a set of size $2^k$.

\begin{example}[A four-bit key]
\label{ch07:fourbit-key}
Let the key space be $\Omega = \set{0,1}^4$, so $N = 2^4 = 16$ keys, from $0000$ to $1111$. Under the uniform distribution each key has probability $1/16$. Eve guesses one key. The event ``Eve guesses right'' has exactly one favorable outcome, so $\Pr[\text{correct}] = 1/16 = 0.0625$. The event ``the key starts with a $1$'' contains the $8$ keys $1000,\dots,1111$, so its probability is $8/16 = 1/2$. We will scale these toy numbers up to $128$ and $256$ bits later, but the arithmetic never changes -- only the exponents do.
\end{example}

We write $k \getsr \set{0,1}^{\ell}$ to mean ``sample $k$ uniformly at random from the set of $\ell$-bit strings.'' That dollar-sign arrow is one of the most important symbols in applied cryptography. Almost every security proof in this book begins by assuming some value was drawn this way, and almost every real-world break begins with the discovery that it was \emph{not}.

\subsection{Conditional probability and Bayes}

Eve's knowledge changes as she observes ciphertexts. The tool for ``probability of $A$ given that I now know $B$'' is conditional probability:
\[
  \Pr[A \mid B] = \frac{\Pr[A \cap B]}{\Pr[B]}, \qquad \text{defined when } \Pr[B] > 0.
\]
Read it as: restrict the universe to the world where $B$ happened, then ask what fraction of that smaller world also has $A$. Rearranging gives the multiplication rule $\Pr[A \cap B] = \Pr[A \mid B]\,\Pr[B]$, and swapping the roles of $A$ and $B$ gives \emph{Bayes' rule}:
\[
  \Pr[A \mid B] = \frac{\Pr[B \mid A]\,\Pr[A]}{\Pr[B]}.
\]
Bayes is how an attacker updates a guess about a key in light of leaked data, and it is also the formal heart of what ``a ciphertext reveals nothing'' should mean: the message should be just as likely after seeing the ciphertext as before. We make that precise when we reach perfect secrecy in Chapter~9; here we just plant the flag.

\begin{example}[A side channel updates a guess]
\label{ch07:bayes-sidechannel}
Return to the four-bit key of Example~\ref{ch07:fourbit-key}, uniform over $16$ values. A noisy power-analysis probe leaks one bit about the key's Hamming weight (the number of $1$s). The probe reports ``the weight is even'' and is correct with probability $0.8$, wrong with probability $0.2$. Of the $16$ keys, exactly $8$ have even weight ($0000, 0011, \ldots, 1111$) and $8$ have odd weight, so before the leak each class has prior probability $1/2$.

Let $A$ be the event ``the true key has even weight'' and $B$ the event ``the probe says even.'' We want $\Pr[A \mid B]$. The likelihoods are $\Pr[B \mid A] = 0.8$ (probe right on an even key) and $\Pr[B \mid \bar A] = 0.2$ (probe wrong on an odd key). The denominator is $\Pr[B] = 0.8 \cdot \tfrac12 + 0.2 \cdot \tfrac12 = 0.5$. Bayes gives
\[
  \Pr[A \mid B] = \frac{0.8 \cdot 0.5}{0.5} = 0.8.
\]
The leak moved Eve from $1/2$ to $0.8$ confidence that the weight is even. Conditioned on that, her best guess targets one of the $8$ even-weight keys, and a single guess there is correct with probability $0.8 \cdot \tfrac18 = 0.1$ -- up from the $1/16 = 0.0625$ she had with no leak. One noisy bit nearly doubled her one-shot odds. This is the engine of every side-channel attack: each leaked measurement is a Bayesian update, and updates compound.
\end{example}

\subsection{Independence}

Two events $A$ and $B$ are \emph{independent} if
\[
  \Pr[A \cap B] = \Pr[A]\cdot \Pr[B],
\]
equivalently $\Pr[A \mid B] = \Pr[A]$: learning $B$ tells you nothing about $A$. Independence is what we demand of the bits coming out of a good random source -- each new bit should be a fresh fair coin, unaffected by all the bits before it. When that assumption fails, keys collapse.

\begin{attackbox}[When ``random'' bits are not independent]
In 2012 two research teams scanned millions of RSA public keys collected from the public Internet and computed pairwise GCDs. Roughly $0.2\%$ of TLS and SSH keys shared a prime factor with some other key, which let anyone factor both and recover the private keys outright. The cause was a broken independence assumption: embedded devices (routers, firewalls) generated key material moments after boot, before their entropy pools had collected enough unpredictable events. Different devices ended up drawing from nearly the same impoverished state, so their ``random'' primes were correlated rather than independent. The math was fine; the probability model the implementers assumed -- fresh independent randomness -- did not hold in the field. We revisit entropy starvation in Chapter~17.
\end{attackbox}

A handy consequence of independence: if you draw $k_1, k_2, \dots, k_n$ as independent uniform bits, then the whole string $k_1 k_2 \cdots k_n$ is uniform over $\set{0,1}^n$. Independence is precisely what turns ``$n$ good coins'' into ``one uniform $n$-bit key.'' Lose independence and you lose the size of your key space, even if every individual bit still looks balanced.

\section{Random variables and expectation}

A \emph{random variable} $X$ is a function from the sample space to numbers, $X : \Omega \to \R$. If you roll a die, ``the number showing'' is a random variable; ``$1$ if the roll is even, else $0$'' is another. The \emph{expectation} (mean) is the probability-weighted average of its values:
\[
  \mathbb{E}[X] = \sum_{\omega \in \Omega} \Pr[\omega]\, X(\omega) = \sum_{x} x \cdot \Pr[X = x].
\]
Expectation is \emph{linear}: $\mathbb{E}[X + Y] = \mathbb{E}[X] + \mathbb{E}[Y]$ always, with no independence required. That one fact makes a surprising number of counting arguments trivial, and we use it below to count expected collisions.

\begin{example}[Expected number of guesses]
\label{ch07:expected-guesses}
Eve attacks the $4$-bit key of Example~\ref{ch07:fourbit-key} by trying keys one at a time in a random order until she hits the right one. Let $X$ be the number of guesses she needs. Each position $1,2,\dots,16$ is equally likely to hold the true key, so $\Pr[X = i] = 1/16$ for each $i$, and
\[
  \mathbb{E}[X] = \sum_{i=1}^{16} i \cdot \tfrac{1}{16} = \tfrac{1}{16}\cdot\tfrac{16\cdot 17}{2} = \tfrac{17}{2} = 8.5.
\]
On average she needs about half the key space, which is the familiar ``brute force costs $2^{\ell-1}$ on average'' rule. For a $128$-bit key that is $2^{127}$ expected guesses, which is why brute force is hopeless.
\end{example}

\subsection{The indicator trick}

A random variable that is $1$ when some event happens and $0$ otherwise is called an \emph{indicator}, written $\mathbf{1}_A$. Its expectation is just the probability of the event: $\mathbb{E}[\mathbf{1}_A] = 1\cdot\Pr[A] + 0\cdot\Pr[\bar A] = \Pr[A]$. Combine that with linearity of expectation and you can count the expected number of times \emph{any} of a pile of events fires, even when the events overlap or are correlated. This is the engine behind the birthday bound, so keep it warm.

\section{Entropy: measuring unpredictability on average}

We need a single number that says ``how hard is this secret to guess.'' Claude Shannon gave us one in 1948, and it remains the right tool for reasoning about average-case unpredictability.

\begin{historybox}
Shannon's 1948 paper ``A Mathematical Theory of Communication'' founded information theory in one stroke, and his companion 1949 paper ``Communication Theory of Secrecy Systems'' did the same for modern cryptography. He proved that the one-time pad is unbreakable and -- more sobering -- that any unbreakable cipher needs a key at least as long as the message. Both results are theorems about entropy. The word ``entropy'' he borrowed from thermodynamics on the advice of John von Neumann, who reportedly told him it would win every argument because nobody really knows what entropy is.
\end{historybox}

\subsection{The definition and where it comes from}

Let $X$ be a random variable taking value $x$ with probability $p_x = \Pr[X = x]$. The \emph{Shannon entropy} of $X$, measured in bits, is
\[
  H(X) = -\sum_{x} p_x \log_2 p_x = \sum_x p_x \log_2 \frac{1}{p_x},
\]
with the convention $0\log_2 0 = 0$ (an impossible outcome contributes nothing). The quantity $\log_2(1/p_x)$ is the \emph{surprise} of outcome $x$: rare outcomes ($p_x$ small) are very surprising (large $\log_2(1/p_x)$), certain outcomes ($p_x = 1$) carry zero surprise. Entropy is the \emph{average surprise} you experience when you learn the value of $X$.

\begin{intuitionbox}
Think of entropy as the average number of yes/no questions an optimal interrogator needs to pin down $X$. If $X$ is one of $8$ equally likely values, a smart questioner splits the candidates in half each time -- ``is it in the top half?'' -- and nails it in $3$ questions: $H(X) = \log_2 8 = 3$ bits. If one value is overwhelmingly likely, the questioner usually guesses it first and is done in roughly one question, so entropy is low. Entropy counts questions; questions are bits.
\end{intuitionbox}

\subsection{Worked entropies on tiny distributions}

\begin{example}[Fair versus biased coin]
\label{ch07:coin-entropy}
A fair coin has $p_{\text{H}} = p_{\text{T}} = 1/2$:
\[
  H = -\tfrac12 \log_2 \tfrac12 - \tfrac12 \log_2 \tfrac12 = \tfrac12 + \tfrac12 = 1 \text{ bit}.
\]
Exactly one bit, as it must be. Now bias the coin so heads comes up with probability $0.9$:
\[
  H = -0.9\log_2 0.9 - 0.1\log_2 0.1 \approx 0.9(0.152) + 0.1(3.32) \approx 0.137 + 0.332 = 0.469 \text{ bits}.
\]
The biased coin still produces ``a bit'' of data per flip, but it carries less than half a bit of genuine unpredictability. If you built a key out of such coins, an attacker who knows the bias would guess each bit right $90\%$ of the time. This gap between raw bit-length and real entropy is where countless ``random'' tokens have died.
\end{example}

\begin{example}[A loaded four-sided die]
\label{ch07:loaded-die}
Let $X$ take four values with probabilities $\tfrac12, \tfrac14, \tfrac18, \tfrac18$. Then
\[
  H(X) = \tfrac12\log_2 2 + \tfrac14\log_2 4 + \tfrac18\log_2 8 + \tfrac18\log_2 8
       = \tfrac12(1) + \tfrac14(2) + \tfrac18(3) + \tfrac18(3) = \tfrac{7}{4} = 1.75 \text{ bits}.
\]
A uniform four-sided die would have entropy $\log_2 4 = 2$ bits, so the loading costs us a quarter-bit. Notice how the most likely outcome ($p = 1/2$) contributes its small surprise of $1$ bit, while the rare outcomes contribute more surprise each but are down-weighted by how seldom they occur.
\end{example}

\subsection{The properties that matter for crypto}

Three facts about Shannon entropy carry their weight throughout the book.

\begin{theorem}[Entropy is maximized by the uniform distribution]
\label{ch07:max-entropy}
For any random variable $X$ over a set of size $N$, $0 \le H(X) \le \log_2 N$. The lower bound is met exactly when $X$ is constant (no uncertainty), and the upper bound is met exactly when $X$ is uniform.
\end{theorem}

\begin{whybox}
Why uniform is the unpredictability champion. Suppose two outcomes had unequal probabilities, say $p$ and $q$ with $p > q$. Nudge a sliver of probability from the likelier one to the rarer one, making them more equal. A short calculus check on $-p\log_2 p - q\log_2 q$ shows this \emph{increases} the total -- surprise grows when you flatten the peaks. The only configuration with no peak left to flatten is the flat one, where every outcome has probability $1/N$. So among all distributions on $N$ values, uniform is the most uncertain, and its entropy is $\log_2 N$. This is exactly why we want keys uniform: any deviation strictly lowers entropy and hands the attacker an edge, however small.
\end{whybox}

The second fact: entropy of independent sources \emph{adds}. If $X$ and $Y$ are independent, $H(X, Y) = H(X) + H(Y)$. Sixteen independent fair coins give $16$ bits of entropy, which is why ``$\ell$ independent uniform bits'' equals ``$\ell$ bits of entropy.'' The third fact, the data-processing inequality: no deterministic function can manufacture entropy, $H(f(X)) \le H(X)$. You cannot post-process your way to more unpredictability than you started with -- a hash of a low-entropy password is still low-entropy, a point attackers exploit ruthlessly in password cracking.

\begin{figure}[H]
\centering
\begin{tikzpicture}[font=\small,>=stealth,scale=0.92]
  % Panel A: General System with Information Leakage
  \begin{scope}[shift={(0,0)}]
    \draw[thick,gray!40,fill=gray!5,rounded corners=6pt] (-2.2,-2.2) rectangle (4.2,2.2);
    \node[font=\bfseries\footnotesize,sourceblue!90!black] at (1.0,1.8) {General Channel: Partial Leakage};

    % Left circle: H(M)
    \draw[thick,sourceblue!80,fill=sourceblue!15] (0,0) circle (1.4cm);
    % Right circle: H(C)
    \draw[thick,targetgreen!80,fill=targetgreen!15] (2.0,0) circle (1.4cm);
    % Intersection: I(M;C)
    \begin{scope}
      \clip (0,0) circle (1.4cm);
      \fill[structurepurple!35] (2.0,0) circle (1.4cm);
    \end{scope}

    \node[font=\scriptsize\bfseries,sourceblue!90!black] at (-0.6,0) {$H(M|C)$};
    \node[font=\scriptsize\bfseries,structurepurple!90!black] at (1.0,0) {$I(M;C)$};
    \node[font=\scriptsize\bfseries,targetgreen!90!black] at (2.6,0) {$H(C|M)$};

    \node[font=\tiny,align=center] at (1.0,-1.6) {
      Mutual info $I(M;C) > 0$\\
      Ciphertext leaks information about $M$.
    };
  \end{scope}

  % Panel B: Perfect Secrecy (Shannon's Condition)
  \begin{scope}[shift={(7.2,0)}]
    \draw[thick,targetgreen!40,fill=targetgreen!5,rounded corners=6pt] (-1.8,-2.2) rectangle (4.6,2.2);
    \node[font=\bfseries\footnotesize,targetgreen!90!black] at (1.4,1.8) {Perfect Secrecy: Complete Decoupling};

    % Left circle: H(M)
    \draw[thick,sourceblue!80,fill=sourceblue!15] (0,0) circle (1.3cm);
    \node[font=\scriptsize\bfseries,sourceblue!90!black,align=center] at (0,0.15) {$H(M)$\\[2pt]\tiny\normalfont Plaintext entropy};

    % Right circle: H(C)
    \draw[thick,targetgreen!80,fill=targetgreen!15] (2.8,0) circle (1.3cm);
    \node[font=\scriptsize\bfseries,targetgreen!90!black,align=center] at (2.8,0.15) {$H(C)$\\[2pt]\tiny\normalfont Ciphertext entropy};

    \node[draw=targetgreen!80,fill=white,rounded corners=2pt,font=\tiny\bfseries,align=center,inner sep=3pt] at (1.4,-0.65) {
      $I(M; C) = 0$\\
      $H(M|C) = H(M)$
    };

    \node[font=\tiny,targetgreen!90!black,align=center] at (1.4,-1.6) {
      \textbf{Zero mutual information:}\\
      Observing $C$ leaves uncertainty $H(M)$ unchanged.
    };
  \end{scope}
\end{tikzpicture}
\caption{Information-theoretic view of cryptographic secrecy. Left: In a leaky or classical cipher, observing ciphertext $C$ reduces uncertainty about message $M$ by the mutual information $I(M;C) > 0$. Right: Under Shannon's condition of perfect secrecy, $M$ and $C$ are statistically independent ($I(M;C) = 0$), so $H(M|C) = H(M)$ and seeing the ciphertext provides zero information about the plaintext.}
\label{ch07:fig-entropy-venn}
\end{figure}

\begin{connectionbox}
Entropy, expectation, and the logarithm all live in the companion probability volume. The $\log_2(1/p)$ ``surprise'' quantity is the information content of an event; averaging it against the distribution is just an expectation, $H(X) = \mathbb{E}[\log_2(1/p_X)]$. If the manipulations in the worked examples felt fast, the chapters on discrete random variables and expectation there walk the same algebra at a slower pace, including the calculus argument behind Theorem~\ref{ch07:max-entropy}.
\end{connectionbox}

\section{Min-entropy: the paranoid measure}

Shannon entropy is an \emph{average}, and averages lie to you in exactly the situation cryptographers care about: the attacker who makes the single best guess. Suppose a key is $0$ with probability $1/2$ and otherwise uniform over the remaining $2^{128} - 1$ values. Its Shannon entropy is enormous, $65$ bits (the $1/2$ spike contributes a half-bit, and the flat tail contributes $\tfrac12\log_2\!\big(2(2^{128}-1)\big) \approx 64.5$). But an attacker who simply guesses ``$0$'' wins half the time on the first try. Shannon entropy says this key is strong; reality says it is catastrophic. We need a measure that reports the worst case, not the average.

\begin{definition}[Min-entropy]
\label{ch07:min-entropy-def}
The \emph{min-entropy} of $X$ is
\[
  H_\infty(X) = -\log_2 \Big( \max_x \Pr[X = x] \Big) = \min_x \log_2 \frac{1}{p_x}.
\]
\end{definition}

In words: find the single most likely outcome, call its probability $p_{\max}$, and report $-\log_2 p_{\max}$. Because the best single guess succeeds with probability exactly $p_{\max} = 2^{-H_\infty(X)}$, min-entropy directly bounds the attacker's one-shot success:
\[
  \Pr[\text{best single guess is correct}] = 2^{-H_\infty(X)}.
\]
That clean operational meaning is why min-entropy, not Shannon entropy, is the quantity randomness extractors and key-derivation standards are written against.

\begin{example}[The two measures disagree]
\label{ch07:hinf-vs-h}
Take the loaded die of Example~\ref{ch07:loaded-die} with probabilities $\tfrac12, \tfrac14, \tfrac18, \tfrac18$. Its Shannon entropy was $1.75$ bits. Its min-entropy uses only the peak, $p_{\max} = 1/2$:
\[
  H_\infty(X) = -\log_2 \tfrac12 = 1 \text{ bit}.
\]
A full three-quarters of a bit lower. The most likely outcome happens half the time, so a one-guess attacker wins $50\%$ of the time, exactly as $2^{-1}$ predicts. For a uniform distribution the two measures coincide -- when every $p_x = 1/N$, both equal $\log_2 N$ -- which is one more reason uniform keys are the gold standard. The measures only diverge when the distribution has a peak, and the peak is precisely what an attacker aims at (Figure~\ref{ch07:loaded-die-bars}).
\end{example}

\begin{figure}[H]
\centering
\begin{tikzpicture}
\begin{axis}[
  width=10cm, height=5cm,
  title style={font=\small},
  title={Outcome Probabilities vs.\ Min-Entropy Peak},
  ybar, bar width=22pt,
  ymin=0, ymax=0.6,
  xtick={1,2,3,4}, xticklabels={$x_1$,$x_2$,$x_3$,$x_4$},
  ylabel={$p_x$}, xlabel={outcome},
  axis lines=left, ymajorgrids, grid style={gray!20},
  nodes near coords, nodes near coords style={font=\small},
  point meta=explicit symbolic,
]
\addplot[fill=trapred!70, draw=trapred] coordinates {(1,0.5) [$\tfrac12$]};
\addplot[fill=sourceblue!50, draw=sourceblue] coordinates {(2,0.25) [$\tfrac14$] (3,0.125) [$\tfrac18$] (4,0.125) [$\tfrac18$]};
\end{axis}
\node[trapred, font=\small\bfseries] at (4.1,3.6) {peak $p_{\max}=\tfrac12$};
\draw[trapred,->,>=Stealth] (3.6,3.4) -- (1.55,3.1);
\end{tikzpicture}
\caption{The loaded die of Example~\ref{ch07:hinf-vs-h}. Shannon entropy ($1.75$ bits) averages the surprise across all four bars; min-entropy ($1$ bit) looks only at the tallest bar, the one an attacker guesses first.}
\label{ch07:loaded-die-bars}
\end{figure}

\begin{formalbox}
In general $0 \le H_\infty(X) \le H(X) \le \log_2 N$. Min-entropy is always the most pessimistic of the entropy family (it is the limit of the Renyi entropies $H_\alpha$ as $\alpha \to \infty$, hence the $\infty$ subscript). When someone in cryptography says ``this source has $k$ bits of entropy'' and means it as a security claim, they should mean min-entropy. If a standard says a seed needs ``$256$ bits of entropy,'' read it as min-entropy unless told otherwise -- that is the reading that makes the security argument go through.
\end{formalbox}

\begin{applicationbox}
NIST Special Publication 800-90B, the U.S. standard for entropy sources feeding random-number generators, is built entirely around min-entropy. A hardware noise source must be \emph{estimated} for its min-entropy, and the estimate must survive a battery of statistical tests designed to catch the source overstating itself. Then a \emph{conditioning} step (often based on a hash or block cipher) compresses many weak samples into fewer, fuller-entropy bits. The reason the standard fixates on min-entropy rather than Shannon entropy is the one-shot-guess bound above: a generator that is usually unpredictable but occasionally spits out a likely value is exactly the failure mode that catches fire in deployment, and only min-entropy sees it.
\end{applicationbox}

\begin{trapbox}
A pile of bytes is not a pile of entropy. A megabyte of the digits of $\pi$ has eight million bits of length and essentially zero entropy -- it is fully determined. Length is what you store; entropy is what the attacker has to guess. Always ask of a ``random'' value: what distribution did it come from, and what is the \emph{largest} probability any single value gets? That second question is the min-entropy question, and it is the one that keeps keys alive.
\end{trapbox}

\subsection{Guessing entropy: how many tries on average}

Min-entropy answers ``how often does the \emph{first} guess win?'' Sometimes you want the dual question: ``if the attacker guesses values in the smartest order, how many guesses until success on average?'' That number is the \emph{guessing entropy} (or guesswork), and it is just an expectation in disguise.

Order the outcomes from most to least likely, $p_1 \ge p_2 \ge \cdots \ge p_M$. An optimal attacker tries them in exactly that order, so the value with probability $p_i$ is found on guess number $i$. The expected number of guesses is
\[
  G(X) = \sum_{i=1}^{M} i \cdot p_i.
\]

\begin{example}[Three ways to score the loaded die]
\label{ch07:guesswork}
Take yet again the die with sorted probabilities $\tfrac12, \tfrac14, \tfrac18, \tfrac18$. Its guessing entropy is
\[
  G(X) = 1\cdot\tfrac12 + 2\cdot\tfrac14 + 3\cdot\tfrac18 + 4\cdot\tfrac18 = \tfrac12 + \tfrac12 + \tfrac38 + \tfrac12 = \tfrac{15}{8} = 1.875.
\]
A uniform draw over $4$ values would need $G = (1+2+3+4)/4 = 2.5$ guesses on average. So we now have three numbers describing the same distribution: $H = 1.75$ bits (average surprise), $H_\infty = 1$ bit (worst-case one-shot), and $G = 1.875$ guesses (average optimal guesswork). They measure genuinely different things, and a security claim has to name which one it means.
\end{example}

\begin{whybox}
Why guesswork is not just $2^{H}$. A tempting shortcut says ``$\ell$ bits of entropy means about $2^\ell$ guesses.'' That is right for the uniform distribution and badly wrong for a peaked one. Massey proved a lower bound $G(X) \ge \tfrac14 \cdot 2^{H(X)} + 1$ when $H(X) \ge 2$, but no matching upper bound exists: you can build distributions with large Shannon entropy whose \emph{first} guess still wins almost always, dragging the average guesswork far below $2^{H}$. The biased-key Example~\ref{ch07:hinf-vs-h} pattern -- one fat spike on a sea of tiny probabilities -- is exactly such a case. When you care about the attacker who batches many guesses (an offline password cracker), guesswork is the right scoreboard; when you care about the single best shot (a one-time token), min-entropy is. Reach for the measure whose operational meaning matches the threat.
\end{whybox}

\section{The birthday paradox and collision security}

Here is a fact that feels wrong the first time you meet it. In a room of just $23$ people, the chance that two of them share a birthday is already better than even. Most people guess you would need closer to $183$ (half of $365$). The intuition fails because we instinctively count the wrong thing -- people instead of \emph{pairs}. This ``paradox'' is not a curiosity; it is the single most important counting argument in symmetric cryptography, because it fixes the security level of every hash function and every randomly chosen nonce.

\subsection{Why 23 is enough}

The trick is to compute the probability of \emph{no} shared birthday and subtract from one. With $n$ people and $N = 365$ equally likely birthdays, line them up and place them one at a time. The first person can have any birthday. The second must dodge the first: probability $\frac{364}{365}$. The third must dodge two: $\frac{363}{365}$. In general the $k$-th must dodge $k-1$ taken days. The probability all $n$ are distinct is the product
\[
  \Pr[\text{all distinct}] = \prod_{k=1}^{n-1} \Big(1 - \frac{k}{N}\Big),
\]
so $\Pr[\text{some collision}] = 1 - \prod_{k=1}^{n-1}\big(1 - \frac{k}{N}\big)$. Plugging in $N = 365$:

\begin{center}
\begin{tabular}{rl}
\toprule
people $n$ & $\Pr[\text{some shared birthday}]$ \\
\midrule
$10$ & $0.117$ \\
$20$ & $0.411$ \\
$23$ & $0.507$ \\
$30$ & $0.706$ \\
$50$ & $0.970$ \\
$70$ & $0.999$ \\
\bottomrule
\end{tabular}
\end{center}

At $n = 23$ we cross one-half. The reason is the number of \emph{pairs}, not people: with $23$ people there are $\binom{23}{2} = 253$ pairs, each a chance for a match, and $253$ is already comfortably more than half of $365$. We count pairs, and pairs grow like $n^2$.

\begin{figure}[H]
\centering
\begin{tikzpicture}
\begin{axis}[
  width=11cm, height=6.2cm,
  title style={font=\small},
  title={Birthday Collision Probability ($N=365$)},
  xlabel={number of people $n$}, ylabel={$\Pr[\text{shared birthday}]$},
  xmin=0, xmax=80, ymin=0, ymax=1.02,
  domain=0:80, samples=120,
  axis lines=left, grid=major, grid style={gray!20},
  every axis plot/.append style={very thick},
  legend pos=south east, legend style={font=\small},
]
\addplot[sourceblue, smooth] {1 - exp(-x*x/730)};
\addlegendentry{$1 - e^{-n^2/(2N)}$, $N=365$}
\addplot[trapred, dashed, samples=2, domain=0:80] {0.5};
\addlegendentry{$\tfrac12$ crossover}
\draw[gray, dotted] (axis cs:23,0) -- (axis cs:23,0.5);
\node[trapred, font=\small] at (axis cs:35,0.42) {$n \approx 23$};
\end{axis}
\end{tikzpicture}
\caption{The collision probability climbs steeply and crosses $\tfrac12$ near $n = 23$. The curve is the $1 - e^{-n^2/(2N)}$ approximation; the steep rise is the $n^2$ growth of the pair count at work.}
\label{ch07:birthday-curve}
\end{figure}

\subsection{The clean bound}

To see the $n^2$ scaling, use the inequality $1 - x \le e^{-x}$ (the exponential always sits above its tangent line at zero). Then
\[
  \Pr[\text{all distinct}] = \prod_{k=1}^{n-1}\Big(1 - \frac{k}{N}\Big)
  \le \prod_{k=1}^{n-1} e^{-k/N}
  = \exp\!\Big(-\frac{1}{N}\sum_{k=1}^{n-1} k\Big)
  = \exp\!\Big(-\frac{n(n-1)}{2N}\Big).
\]
Therefore the collision probability is at least
\[
  \Pr[\text{collision}] \;\ge\; 1 - \exp\!\Big(-\frac{n(n-1)}{2N}\Big) \;\approx\; 1 - \exp\!\Big(-\frac{n^2}{2N}\Big).
\]
Set the exponent to about $1$ to get a fifty-fifty collision: that happens when $n^2 \approx 2N$, i.e. when
\[
  \boxed{\,n \approx \sqrt{2N} \approx 1.177\sqrt{N}.\,}
\]
For birthdays, $\sqrt{2 \cdot 365} \approx 27$, and the more careful constant $1.177\sqrt{365} \approx 22.5$ -- right on top of the true crossover at $23$ (Figure~\ref{ch07:birthday-curve}). The headline is the square root: you need on the order of $\sqrt{N}$ samples, not $N$, before a collision is likely.

\begin{whybox}
The square root is the whole story for crypto. If outputs are drawn uniformly from a space of size $N = 2^n$ (an $n$-bit hash), then collisions become likely after roughly $\sqrt{2^n} = 2^{n/2}$ samples. The $n$-bit space gives you only $n/2$ bits of \emph{collision} resistance. This is not a weakness you can engineer away with a cleverer hash; it is a consequence of counting pairs, and it applies to every function with a fixed-size output. It is the reason hash output lengths are chosen to be \emph{twice} the desired security level.
\end{whybox}

\subsection{Counting collisions with the indicator trick}

There is a second, even quicker way to land on $\sqrt{N}$ that shows off the indicator trick from earlier. Throw $n$ samples into $N$ uniform bins. For each unordered pair $\set{i, j}$ of samples, let $\mathbf{1}_{ij}$ be the indicator that samples $i$ and $j$ collide. Two uniform draws land in the same bin with probability $1/N$, so $\mathbb{E}[\mathbf{1}_{ij}] = 1/N$. By linearity of expectation -- no independence needed, even though the pair-events overlap -- the expected number of colliding pairs is
\[
  \mathbb{E}[\text{collisions}] = \sum_{\set{i,j}} \mathbb{E}[\mathbf{1}_{ij}] = \binom{n}{2}\cdot\frac{1}{N} = \frac{n(n-1)}{2N}.
\]
This expected count passes $1$ exactly when $n(n-1) \approx 2N$, i.e. $n \approx \sqrt{2N}$ again. Same threshold, derived in two lines. When the expected number of collisions crosses one, collisions stop being rare.

\subsection{What this dictates for hash functions}

A cryptographic hash $\Hash : \set{0,1}^* \to \set{0,1}^n$ should be \emph{collision resistant}: nobody should find $x \ne y$ with $\Hash(x) = \Hash(y)$. The birthday bound says a generic attacker -- one who just hashes random inputs and waits for a repeat, knowing nothing about the hash's internals -- finds a collision in about $2^{n/2}$ work. So:

\begin{itemize}
  \item A $128$-bit hash gives only $\approx 2^{64}$ collision resistance. $2^{64}$ operations is large but has been within reach of well-funded effort for years. Treat $128$-bit digests as obsolete for collision resistance.
  \item SHA-$256$ outputs $256$ bits, for $\approx 2^{128}$ collision resistance -- the modern comfortable target.
  \item If you need $256$-bit security against collisions (e.g. long-term or post-quantum margins), reach for a $512$-bit digest like SHA-$512$.
\end{itemize}

Note the asymmetry that trips people up. \emph{Preimage} resistance (given a digest, find any input mapping to it) costs $\approx 2^n$ -- the full space, because you are aiming at one fixed target, not waiting for any pair to coincide. \emph{Collision} resistance costs only $\approx 2^{n/2}$. The birthday bound bites collisions, not preimages, which is why a hash is always weaker against collisions than its length naively suggests.

\begin{attackbox}[Birthday attacks that actually happened]
The birthday bound is not theoretical. MD5 produces a $128$-bit digest, so its generic collision resistance is only $2^{64}$ -- and structural weaknesses dropped real attacks far below even that. In 2008 researchers used MD5 collisions to forge a rogue certificate-authority certificate, letting them mint trusted certificates for any website. The 2012 Flame espionage malware used an MD5 collision to forge a Microsoft code-signing certificate and spread via Windows Update. SHA-1, with a $160$-bit digest ($2^{80}$ generic), fell in 2017 to the \emph{SHAttered} attack, which produced two different PDF files with the same SHA-1 hash; in 2020 a \emph{chosen-prefix} SHA-1 collision broke real PGP trust. In every case the underlying clock was the birthday bound: a digest only half as strong as its length suggests, eroded further by cryptanalysis.
\end{attackbox}

\begin{connectionbox}
The companion probability book derives the $\sqrt{N}$ scaling from first principles and proves the $1 - e^{-x}$ bound used here, and the algorithms book pairs it with the \emph{rho} method (Pollard's cycle-finding) that turns the birthday bound into a memory-cheap $O(2^{n/2})$ collision-finding algorithm -- the practical engine behind the attacks above. If you want the variance of the collision count and tail bounds (how concentrated the ``about $2^{n/2}$'' really is), that is the second-moment material there.
\end{connectionbox}

\begin{implementbox}
A tiny birthday-collision experiment. We hash random inputs with a deliberately short $24$-bit digest and watch how many tries it takes to collide, then compare with the prediction $\approx 2^{n/2} = 2^{12} = 4096$. This is illustration code for learning, never for production -- truncating SHA-256 to $24$ bits is exactly the kind of thing you must never ship.

\begin{lstlisting}[language=Python]
import hashlib, os

def short_hash(data, nbits=24):
    full = hashlib.sha256(data).digest()
    val = int.from_bytes(full, "big")
    return val & ((1 << nbits) - 1)      # keep only the low nbits

def find_collision(nbits=24):
    seen = {}                            # digest -> input that produced it
    tries = 0
    while True:
        x = os.urandom(16)
        h = short_hash(x, nbits)
        tries += 1
        if h in seen and seen[h] != x:
            return tries, seen[h], x      # collision found
        seen[h] = x

n = 24
tries, a, b = find_collision(n)
print(f"collided after {tries} hashes; birthday estimate ~ {2**(n//2)}")
# Typical run: a few thousand hashes, right around 2**12 = 4096.
\end{lstlisting}
Run it a dozen times and average: the collision count clusters near $2^{12}$, not near $2^{24}$. You are watching the square root with your own eyes.
\end{implementbox}

\begin{applicationbox}
Birthday math sets concrete limits in protocols you use daily. AES in GCM mode uses a $96$-bit nonce; a \emph{random} nonce risks repetition by the birthday bound after about $2^{48}$ messages under one key, and a repeated GCM nonce is catastrophic -- it leaks the authentication key. That is why TLS 1.3 derives GCM nonces from a counter rather than gambling on randomness. The older $64$-bit block ciphers (DES, 3DES, Blowfish) suffer the \emph{Sweet32} attack: in CBC mode, ciphertext blocks start colliding after $\approx 2^{32}$ blocks ($\approx 32$ GB), and each collision leaks a relationship between plaintexts. A long-lived HTTPS or VPN session over 3DES could be partially decrypted purely from birthday collisions, which is why $64$-bit block ciphers were retired from TLS.
\end{applicationbox}

\section{Statistical distance and indistinguishability}

The final tool is the formal way to say two distributions ``look the same.'' This is the backbone of nearly every security definition in modern cryptography: a scheme is secure when the distribution an attacker sees in the real world is indistinguishable from some ideal, secret-free distribution.

\subsection{Definition and the guessing-game meaning}

\begin{definition}[Statistical distance]
\label{ch07:stat-distance}
Let $P$ and $Q$ be distributions over the same finite set $\Omega$. Their \emph{statistical distance} (also called total variation distance) is
\[
  \Delta(P, Q) = \frac{1}{2}\sum_{\omega \in \Omega} \abs{P(\omega) - Q(\omega)} = \max_{A \subseteq \Omega} \big( P(A) - Q(A) \big).
\]
\end{definition}

The two forms are equal, and each tells a story. The sum form adds up, over every outcome, how much the two distributions disagree, halved so the total never exceeds $1$. The max form is the operational one: it says $\Delta(P,Q)$ is the \emph{best possible advantage} any test can achieve at telling the two distributions apart. $\Delta = 0$ means the distributions are identical and no test does better than a coin flip; $\Delta = 1$ means their supports are disjoint and one sample reveals which is which.

\begin{intuitionbox}
Imagine a game. A referee secretly flips a fair coin: heads, she hands you a sample from $P$; tails, a sample from $Q$. You inspect the sample and shout your guess. If $P$ and $Q$ are identical you are reduced to guessing, winning half the time. The more they differ, the better you can do -- and the very best win probability you can ever achieve is exactly $\tfrac12 + \tfrac12\Delta(P,Q)$. Statistical distance is your \emph{edge over random guessing}, doubled. That is what ``$0.01$ statistical distance'' buys an attacker: a $1\%$ edge, and not a sliver more.
\end{intuitionbox}

\begin{figure}[H]
\centering
\begin{tikzpicture}[font=\small,>=stealth,scale=0.95]
  % Referee
  \node[draw=structurepurple!80,fill=structurepurple!12,thick,rounded corners=4pt,minimum width=2.4cm,minimum height=1.3cm,align=center] (ref) at (0,0) {
    \textbf{Referee}\\
    \scriptsize Flips coin $b \getsr \{0,1\}$
  };

  % Worlds P and Q
  \node[draw=sourceblue!80,fill=sourceblue!12,thick,rounded corners=3pt,minimum width=2.6cm,minimum height=0.9cm,align=center] (P) at (4.5,1.1) {
    \textbf{World 0:}\\
    Sample $x \getsr P$
  };
  \node[draw=targetgreen!80,fill=targetgreen!12,thick,rounded corners=3pt,minimum width=2.6cm,minimum height=0.9cm,align=center] (Q) at (4.5,-1.1) {
    \textbf{World 1:}\\
    Sample $x \getsr Q$
  };

  % Distinguisher
  \node[draw=trapred!80,fill=trapred!12,thick,rounded corners=4pt,minimum width=2.6cm,minimum height=1.3cm,align=center] (D) at (9.0,0) {
    \textbf{Distinguisher $D$}\\
    \scriptsize Inspects sample $x$\\
    \scriptsize Outputs guess $\hat b \in \{0,1\}$
  };

  % Arrows
  \draw[->,very thick,sourceblue] (ref.north east) to[out=30,in=180] node[above,font=\tiny\bfseries]{if $b=0$} (P.west);
  \draw[->,very thick,targetgreen] (ref.south east) to[out=-30,in=180] node[below,font=\tiny\bfseries]{if $b=1$} (Q.west);

  \draw[->,very thick,sourceblue] (P.east) to[out=0,in=150] node[above,font=\tiny\bfseries]{$x \sim P$} (D.north west);
  \draw[->,very thick,targetgreen] (Q.east) to[out=0,in=-150] node[below,font=\tiny\bfseries]{$x \sim Q$} (D.south west);

  % Bottom banner
  \node[draw=gray!50,fill=gray!5,rounded corners=3pt,thick,align=center,font=\footnotesize,inner sep=4pt] at (4.5,-2.3) {
    \textbf{Game Outcome:} Distinguisher wins if $\hat b = b$.\\
    $\Pr[\text{Distinguisher wins}] \le \tfrac{1}{2} + \tfrac{1}{2}\Delta(P, Q)$. When $\Delta = 0$, success is capped at $50\%$.
  };
\end{tikzpicture}
\caption{The statistical distinguishing game. A referee flips a fair bit $b$ and provides a sample from either distribution $P$ or $Q$. The distinguisher inspects the sample and guesses which world produced it. The maximum possible advantage over a blind guess is precisely half the statistical distance $\Delta(P,Q)$.}
\label{ch07:referee-game}
\end{figure}

\subsection{A worked distance}

\begin{example}[Two coins]
\label{ch07:two-coins}
Let $P$ be a fair coin, $P(\text{H}) = P(\text{T}) = \tfrac12$, and $Q$ a biased coin, $Q(\text{H}) = \tfrac34, Q(\text{T}) = \tfrac14$. Then
\[
  \Delta(P, Q) = \tfrac12\Big( \abs{\tfrac12 - \tfrac34} + \abs{\tfrac12 - \tfrac14} \Big)
  = \tfrac12\Big( \tfrac14 + \tfrac14 \Big) = \tfrac14.
\]
Check it against the max form: the best distinguisher guesses ``$Q$'' whenever it sees heads. It is right with probability $Q(\text{H}) = \tfrac34$ when the sample truly came from $Q$, and wrong (calls $Q$ on a head) with probability $P(\text{H}) = \tfrac12$ when it came from $P$. The advantage is $\tfrac34 - \tfrac12 = \tfrac14$, matching $\Delta$. Best win probability $\tfrac12 + \tfrac12\cdot\tfrac14 = \tfrac58 = 0.625$ in the referee game of Figure~\ref{ch07:referee-game}. So one look at a single flip lets you beat the coin flip by $12.5$ percentage points -- meaningful, because these distributions are crude. Cryptographic distributions are tuned so $\Delta$ is negligibly small.
\end{example}

\begin{trapbox}
``Small distance per sample'' does not mean ``small distance forever.'' Statistical distance is not free under repetition. If a single sample of $P$ sits at distance $\delta$ from a single sample of $Q$, then $m$ independent samples sit at distance \emph{at most} $m\delta$ -- and often close to it. A pseudorandom generator whose one-block output is distance $2^{-40}$ from uniform is not distance $2^{-40}$ from uniform after a billion blocks; the gap can grow toward $2^{30}\cdot 2^{-40} = 2^{-10}$, a thousand times larger. This is why security proofs track how many queries the attacker makes and why ``rekey before the bound'' rules exist. A distance that is negligible per use can become a real edge once an attacker harvests enough uses.
\end{trapbox}

\subsection{From statistical to computational indistinguishability}

Statistical distance is an \emph{information-theoretic} measure: it bounds \emph{any} distinguisher, even one with unlimited computing power. That is a high bar, and only a few schemes (the one-time pad, secret sharing) clear it. Most of cryptography settles for something weaker but sufficient: distributions that are \emph{computationally} indistinguishable, meaning no \emph{efficient} (polynomial-time) algorithm can tell them apart with non-negligible advantage, even though an all-powerful adversary could.

\begin{formalbox}
We say ensembles $\set{P_\lambda}$ and $\set{Q_\lambda}$ (indexed by a security parameter $\lambda$, e.g. key length) are \emph{computationally indistinguishable}, written $P_\lambda \approx_c Q_\lambda$, if for every $\PPT$ distinguisher $D$,
\[
  \abs{\,\Pr[D(P_\lambda) = 1] - \Pr[D(Q_\lambda) = 1]\,} \le \negl(\lambda),
\]
where $\negl(\lambda)$ shrinks faster than any inverse polynomial. Statistical indistinguishability ($\Delta \le \negl$) is the special case where the bound holds against \emph{all} distinguishers, not just efficient ones. A pseudorandom generator's defining promise is exactly this: its output is computationally indistinguishable from uniform, even though, having a short seed, it is statistically very far from uniform. We develop the computational story in Chapter~8.
\end{formalbox}

\begin{whybox}
Why indistinguishability is the right notion of ``secure.'' We cannot demand that a ciphertext literally be independent of the message -- that forces one-time-pad-length keys (Shannon's theorem). Instead we demand something operationally just as good: the ciphertext distribution when encrypting message $m_0$ should be indistinguishable from the distribution when encrypting $m_1$. If no efficient attacker can tell which message was sent, the encryption leaks nothing an attacker can \emph{use}, even if information is technically present. Every ``IND'' security definition in this book -- IND-CPA, IND-CCA -- is a statistical-distance-style game where the attacker's advantage is required to be negligible. Statistical distance is the warm-up; the bounded-distinguisher version is the real thing.
\end{whybox}

\subsection{The data-processing inequality for distance}

One property of statistical distance is so useful it deserves a name. Post-processing cannot increase statistical distance: for any (possibly randomized) function $f$,
\[
  \Delta\big(f(P), f(Q)\big) \le \Delta(P, Q).
\]
If you cannot tell $P$ from $Q$, you cannot tell apart anything you compute from them either -- mashing the samples through $f$ can only blur the difference, never sharpen it. This is what licenses the hybrid arguments that dominate security proofs: swap one component for an indistinguishable one, and the distance the whole system moves is at most the distance you introduced. Distances along a chain of swaps simply add (the triangle inequality, $\Delta(P,R) \le \Delta(P,Q) + \Delta(Q,R)$), so a proof with a handful of indistinguishable steps keeps the total advantage negligibly small.

\begin{connectionbox}
Statistical distance is the total-variation distance from the probability book's chapter on metrics between distributions, and the triangle inequality and data-processing inequality are proved there in full generality. The ``advantage = edge over guessing'' identity is the Neyman-Pearson lemma in disguise -- the optimal distinguisher thresholds the likelihood ratio $P(\omega)/Q(\omega)$, which is precisely the rule ``guess $Q$ when you see heads'' from Example~\ref{ch07:two-coins}.
\end{connectionbox}

\section{Putting the pieces together}

Every chapter ahead leans on something here. When we build a stream cipher, ``the keystream is indistinguishable from random'' is a statistical-distance claim. When we size a hash, ``$n$-bit digest, $n/2$-bit collision security'' is the birthday bound. When we specify a key generator, ``$256$ bits of min-entropy'' is Definition~\ref{ch07:min-entropy-def}, not Shannon entropy. And when we prove the one-time pad perfectly secret in Chapter~9, the proof is a single line of statistical distance: the ciphertext distribution is identical for every message, so $\Delta = 0$, so no attacker -- bounded or not -- learns anything.

\begin{recoverybox}
Before moving on, make sure you can answer these without notes: (1) Why does a uniform key maximize the attacker's uncertainty? (2) What does min-entropy measure that Shannon entropy misses, and why does an attacker care? (3) Why does an $n$-bit hash give only $n/2$-bit collision resistance? (4) What does statistical distance $0.01$ mean for an attacker trying to distinguish two distributions? If any of these is fuzzy, reread that section -- the rest of the book treats them as reflexes.
\end{recoverybox}

\section{Exercises}

\begin{exercise}[Uniform key arithmetic]
\label{ch07:ex-uniform}
A system uses a $6$-bit key, drawn uniformly. (a) How many keys are there, and what is the probability Eve guesses the key on a single try? (b) What is the expected number of guesses for a brute-force attacker trying keys in random order without repeats? (c) Eve learns that the key is even (last bit $0$). After this leak, what is her single-guess success probability, and how many bits of the key's entropy has she gained?
\end{exercise}

\begin{exercise}[A Bayesian leak]
\label{ch07:ex-bayes}
A $4$-bit key is uniform over $16$ values. A side channel reports the parity of the key (XOR of all four bits) and is correct with probability $0.75$, wrong with probability $0.25$. The probe says ``parity is $0$.'' (a) Using Bayes, compute the posterior probability that the true parity is actually $0$. (b) Conditioned on the report, what is Eve's best single-guess success probability, and how does it compare with her $1/16$ before the leak? (c) A second, independent probe of the same accuracy also reports ``parity is $0$.'' Update again and give the new posterior. What does the trend say about repeated noisy measurements?
\end{exercise}

\begin{exercise}[Three scores for one distribution]
\label{ch07:ex-guesswork}
For the distribution with probabilities $0.5, 0.25, 0.125, 0.125$: (a) compute the Shannon entropy $H$, the min-entropy $H_\infty$, and the guessing entropy $G = \sum_i i\,p_i$ (outcomes sorted by decreasing probability). (b) Compute the same three for the uniform distribution on $4$ values. (c) Explain in two sentences which of the three numbers you would quote to an attacker who gets exactly one guess, and which to an attacker running an offline cracker that tries candidates in the best order.
\end{exercise}

\begin{exercise}[Entropy by hand]
\label{ch07:ex-entropy}
Compute the Shannon entropy $H(X)$ and the min-entropy $H_\infty(X)$ for each distribution. (a) $X$ uniform on $8$ values. (b) Probabilities $\tfrac12, \tfrac14, \tfrac14$. (c) Probabilities $0.7, 0.1, 0.1, 0.1$. In which case is the gap $H - H_\infty$ largest, and explain in one sentence why that distribution is the most dangerous to use as a key despite a respectable Shannon entropy.
\end{exercise}

\begin{exercise}[Min-entropy is what bounds guessing]
\label{ch07:ex-minent-bound}
A password is chosen so that ``$password123$'' has probability $0.05$ and the remaining probability is spread uniformly over $10^{6} - 1$ other strings. (a) Compute $H_\infty$. (b) Compute (approximately) $H$. (c) An attacker gets one guess. What is the best success probability, and which entropy predicted it? (d) Argue in two sentences why quoting the Shannon entropy here would dangerously overstate security.
\end{exercise}

\begin{exercise}[The biased-bit key]
\label{ch07:ex-biased-key}
A flawed generator produces each key bit independently, but each bit is $1$ with probability $0.6$ rather than $0.5$. The key is $128$ bits. (a) What is the Shannon entropy per bit, and the total over $128$ bits? (b) What is the min-entropy per bit (hint: the most likely value of each bit is $1$), and the total? (c) The key has $128$ bits of length but how many effective bits of security against a best-guess attacker? Compare to a correct uniform $128$-bit key.
\end{exercise}

\begin{exercise}[Birthday crossover]
\label{ch07:ex-birthday}
Using the approximation $\Pr[\text{collision}] \approx 1 - e^{-n^2/(2N)}$: (a) For $N = 365$, find the smallest $n$ making the collision probability at least $\tfrac12$, and compare with the exact answer $23$. (b) For a $40$-bit hash ($N = 2^{40}$), about how many random inputs before a collision is more likely than not? (c) State the general rule for an $n$-bit hash in one line.
\end{exercise}

\begin{exercise}[Sizing a hash]
\label{ch07:ex-sizing}
You are designing a system that must resist collision-finding by an attacker capable of $2^{100}$ hash evaluations. (a) What is the minimum digest length you should choose, and why? (b) Your colleague argues a $128$-bit digest is fine ``because $128 > 100$.'' Explain precisely where the reasoning breaks. (c) The same system needs preimage resistance against the same $2^{100}$ attacker -- what digest length does that requirement alone demand, and why is it different?
\end{exercise}

\begin{exercise}[Counting collisions with indicators]
\label{ch07:ex-indicator}
Prove, using indicator random variables and linearity of expectation, that when $n$ items are thrown uniformly and independently into $N$ bins, the expected number of colliding (unordered) pairs is $\binom{n}{2}/N$. Then find the $n$ at which this expectation first reaches $1$ for $N = 2^{64}$, and connect it to the $2^{n/2}$ rule.
\end{exercise}

\begin{exercise}[Statistical distance]
\label{ch07:ex-statdist}
Let $P$ be uniform on $\set{1,2,3,4}$ and let $Q$ assign probabilities $0.4, 0.3, 0.2, 0.1$ to those same four values. (a) Compute $\Delta(P, Q)$ from the sum form. (b) Find the optimal distinguisher and its advantage, and verify it equals $\Delta$. (c) What is the best probability of correctly identifying the source in the referee game of this chapter?
\end{exercise}

\begin{exercise}[Distance cannot grow]
\label{ch07:ex-dpi}
Let $P$ and $Q$ be distributions on $\set{0,1}^{8}$ with $\Delta(P, Q) = 0.2$. Let $f$ map each byte to its parity (XOR of its bits), a single output bit. (a) Argue that $\Delta(f(P), f(Q)) \le 0.2$ without computing anything. (b) Give a concrete $P, Q$ where the inequality is strict (post-processing strictly loses distinguishing power), and one where equality holds. (c) Explain in one sentence why this inequality is what makes hybrid security proofs sound.
\end{exercise}

\begin{exercise}[Why indistinguishability, not independence]
\label{ch07:ex-ind}
Shannon proved that perfect secrecy (ciphertext statistically independent of plaintext) forces the key to be at least as long as the message. (a) State why this makes perfect secrecy impractical for, say, encrypting a $1$ GB file with a short key. (b) Explain how computational indistinguishability sidesteps the limit while still giving meaningful security, and identify which assumption from Chapter~8 the relaxation rests on. (c) In one sentence, relate ``IND-CPA security'' to the statistical-distance guessing game of this chapter.
\end{exercise}
